\documentclass[10pt,a4paper]{amsart}
\usepackage[margin=19mm]{geometry}
\usepackage{amsmath,amssymb,mathtools}
\usepackage[scr=rsfs,cal=euler]{mathalfa}
\usepackage{xfrac}
\usepackage[unicode,hidelinks,bookmarks=false]{hyperref}
\newcommand{\ie}{i.e.\ }
\newcommand{\cf}{cf.\ }
\newcommand{\resp}{respectively\ }
\newcommand{\Subsection}{\subsection}
\providecommand{\nobreakdash}{\nobreak\hbox{-}}
\setlength{\emergencystretch}{2em}
\multlinegap0pt
\usepackage{amssymb}
\usepackage{enumerate}
\usepackage{float}
\usepackage{mathtools}
\usepackage{bm}
\usepackage{xcolor}
\usepackage[shortlabels]{enumitem}
\usepackage{tikz}
\newtheorem{theorem}{Theorem}
\newtheorem{cor}{Corollary}
\newtheorem{lemma}{Lemma}
\newtheorem{claim}{Claim}
\newcommand{\Red}{\operatorname{Red}}
\newcommand{\dem}{\operatorname{Dem}}
\newcommand{\lub}{\operatorname{lub}}
\title{Standard Monomials for Positroid Varieties}
\author{Ayah Almousa, Shiliang Gao, Daoji Huang}
\date{Source: arXiv:2309.15384v3 (CC BY 4.0)}
\begin{document}
\maketitle

\noindent\emph{Source and licence.} Standard Monomials for Positroid Varieties. Version arXiv:2309.15384v3 (CC BY 4.0), distributed by arXiv under the Creative Commons Attribution 4.0 International licence, http://creativecommons.org/licenses/by/4.0/, as recorded in the arXiv licence metadata for this exact version and shown on the abstract page https://arxiv.org/abs/2309.15384v3. Full licence notice: \texttt{licenses/arxiv-2309.15384-CC-BY-4.0.txt}.\par\smallskip

\noindent\textit{Editorial scope.} Counted unit: the article's theorem thm:lub (source0.tex lines 3314--3322). The article's complete original proof of it is delivered with it (source0.tex lines 3362--3413, one \texttt{proof} environment of 52 lines, which the article places away from the statement). No mathematics is written, repaired or re-derived: the statement and the proof are the article's own lines, in the article's own notation, and no retained line is substituted or re-flowed.  Retained uncounted as the source-local context the proof needs: lines 3336--3338; lines 3352--3355.  Nothing was omitted from any retained span, so the package carries no omission record.  No retained line contains a citation, so the wrapper carries no bibliography and no bibliography entry.  No retained line contains a figure environment, an image-inclusion command or drawing code, so no image is delivered and the package declares no asset dependency.  Editorial formatting only, in addition: an AMS \texttt{amsart} preamble that restates the article's own declarations and macros used by the retained lines, namely the environments \texttt{theorem}, \texttt{cor}, \texttt{lemma}, \texttt{claim} on separate counters, together with the standard packages those lines need.  The retained environments print in this document as Theorem 1, Corollary 1, Lemma 1, Claim 1 in order of appearance, whereas the article prints the same items with the numbering of its own full document.  The complete native source of the article as distributed in the arXiv e-print for this exact version is bundled unchanged as \texttt{sources/147611/source0.tex}.
\begin{cor}\label{cor:lift}
    For simple reflections $s,s'$ such that $w\leq sw,ws'$, then either $sw,ws'\leq sws'$ or $w = sws'$. 
\end{cor}

\begin{lemma}
\label{lem:dem}
    Let $Q$ be a not necessarily reduced word that contains a subword $Q_w$ that is a reduced word for $w\in W$. Then $\dem(Q)\ge w$.
\end{lemma}

\begin{theorem}\label{thm:lub}
Let \(W\) be a Coxeter group with Bruhat order \(\le\), and let \(W_P\subseteq W\) be a parabolic subgroup.
Fix \(f\in W\). For \(x,y\in W\), let \(\lub(x,y)\) denote the set of \emph{least} upper bounds of \(\{x,y\}\)
in Bruhat order (possibly empty, possibly containing multiple elements).
If \(f_1,f_2\in W_P f\), then every least upper bound of \(\{f_1,f_2\}\) lies in the same coset:
\[
\lub(f_1,f_2)\subseteq W_P f.
\]
\end{theorem}

\begin{proof}[Proof of Theorem~\ref{thm:lub}]
    It is enough to assume that $f\in W^P$ is a minimal coset representative. Let $w_1,w_2\in W_P$ be elements such that $f_1 = w_1f, f_2 = w_2f$. 
    
    We proceed by induction on $\ell(f)$. The base case is $f=e$, the identity. Assume $w\in \lub(w_1,w_2)$. We show that $w\in W_P$. Suppose to the contrary that $w\not\in W_P$ and assume $Q$ is a reduced word of $w$. Then there exists a simple reflection $s_k\not \in W_P$ which appears in the reduced word of $Q$. Let $Q'$ be the (not necessarily reduced) word obtained from $Q$ by crossing out this $s_k$ and let $w'=\dem(Q')$. By assumption, $Q'$ contains a reduced subword of $w_i$ for $i=1,2$. By Lemma~\ref{lem:dem}, we have $w\ge w_i$ for $i=1,2$. Since $w'<w$, it contradicts that $w\in \lub(w_1, w_2)$.

    Suppose the statement holds for all $f\in W^P$ with $\ell(f)< \ell$. We first consider the case where there exists $s_i \in D_{R}(g)\cap D_{R}(f)$. Notice that we also have $s_i \in D_{R}(w_1f),D_{R}(w_2f)$. By the lifting property, $gs_i\geq w_1fs_i,w_2fs_i$. If $gs_i\notin \lub(w_1fs_i,w_2fs_i)$, we can find $h$ where 
    \[w_1fs_i,w_2fs_i \leq h<gs_i.\]
    Since $s_i\in D_{R}(h\ast s_i)\setminus (D_{R}(w_1fs_i)\cup D_{R}(w_2fs_i))$, by the lifting property, we get 
    \[h\ast s_i \geq w_1f,w_2f.\]
    This contradicts the assumption that $g\in \lub(w_1f,w_2f)$ since $g>h\ast s_i$. Therefore $gs_i\in \lub(w_1fs_i,w_2fs_i)$.  Now by the inductive hypothesis, we have $gs_i = wfs_i$ for some $w\in W_P$ and thus $g = wf$ as desired.

    We are now left with the case where $D_{R}(g) \cap D_{R}(f) = \emptyset$. For any $s_i\in D_{R}(g)$, if $s_i\notin D_{R}(w_1f)\cup D_{R}(w_2f)$, then, by the lifting property, $g>gs_i\geq w_1f,w_2f$. This is impossible since $g\in \lub(w_1f,w_2f)$. As a result, we assume without loss that $s_i\in D_{R}(w_1f)$. Set
    \[f' = \begin{cases}
        w_2f & \text{ if }s_i\notin D_{R}(w_2f)\\
        w_2fs_i & \text{otherwise}
    \end{cases}.\]
    \begin{claim}\label{claim:fsi}
        $fs_i\in W_P f$, hence $w_1fs_i,f'\in W_Pf$.
    \end{claim}
    \begin{proof}
        Pick any $a_1\dots a_{\ell}\in \Red(f)$ and $b_1\dots b_m\in \Red(w_1)$, then since $w_1\in W_P$ and $f\in W^P$, $b_1\dots b_m a_1\dots a_\ell \in \Red(w_1f)$. Since $s_i\notin D_{R}(f)$, $a_1\dots a_\ell\,  i \in \Red(fs_i)$. Since $\ell(w_1fs_i) = \ell(w_1f) - 1$, there is a unique $j\in [m]$ such that 
        \[s_{b_j}\cdots s_{b_m}fs_i< s_{b_{j+1}} \cdots s_{b_m}fs_i.\]
        Set $u = s_{b_{j+1}} \cdots s_{b_m}f\in W$, we have 
        \[u\leq s_{b_j}u, us_i\text{ and }us_i>s_{b_j}us_i.\]
        By Corollary~\ref{cor:lift}, $u = s_{b_j}us_{i}$. Since $s_{b_i}\in W_P$ for all $i\in [m]$, we have $u\in W_P f$ and $s_{b_j}us_i\in W_P fs_i$. 
        Therefore $W_Pf=W_Pfs_i$ and thus 
        $w_1fs_i\in W_P f$.
        If $f' = w_2f$, it is clear that $f'\in W_Pf$. Otherwise, $f'\in W_Pf$ follows from the same reasoning as above.
    \end{proof}
    
    \begin{claim}\label{claim:gsi}
         $gs_i\in \lub(w_1fs_i,f')$. 
    \end{claim}
    \begin{proof}
        We first show that $gs_i\geq w_1fs_i,f'$. Since $s_i\in D_{R}(g)\cap D_{R}(w_1f)$ and $g\geq w_1f$, by the lifting property we have $gs_i\geq w_1fs_i$. Since $f'\leq w_2f\leq g$ and $s_i\in D_{R}(g)\setminus D_{R}(f')$, $f'<g$.
        Again by the lifting property, $gs_i\geq f'$.
        

        Suppose that $gs_i$ is not a least upper bound, and pick $h<gs_i$ such that $h\geq w_1fs_i,f'$. Consider $h' = h\ast s_i$. Since $g>gs_i>h$, by the lifting property we have $g>h'$. Since $s_i\in D_{R}(h')\setminus (D_{R}(w_1fs_i)\cup D_{R}(f'))$, again by the lifting property we have $h'\geq w_1f,w_2f$. This contradicts the assumption that $g\in \lub(w_1f,w_2f)$. Therefore $gs_i\in \lub(w_1fs_i,f')$. 
    \end{proof}
    
    
    
    We can now finish the proof by induction on $\ell(w_1)$ and $\ell(w_2)$. If either $\ell(w_1)$ or $\ell(w_2)$ is $0$, then $g = \max(w_1f,w_2f)$ and the theorem statement holds. By Claim~\ref{claim:fsi}, we write $$w_1fs_i = w_1'f\text{ and }f' = w_2'f$$
    where 
    $$w_1',w_2'\in W_P\text{ and }\ell(w_1')<\ell(w_1),\ell(w_2')\leq \ell(w_2).$$ 
    By Claim~\ref{claim:gsi}, $gs_i\in \lub(w_1'f,w_2'f)$ and by the inductive hypothesis, 
    $gs_i = w'f$ for some $w'\in W_P$. We are then done by Claim~\ref{claim:fsi} since
    \[
g = gs_is_i = w'fs_i \in W_Pf.\qedhere
    \]
\end{proof}

\end{document}
